\section{概览}
\begin{enumerate}
  \item 分数占比最高。
  \item  备考方法。
        \begin{enumerate}
          \item 体系化。
          \item 图示化。\zd{民商事法律关系一般都在 3 个以上当事人，得把两两当事人的法律关系搞清楚}。
          \item 案例化。
          \item 只要跟你正常的朴素的法律常识不认识一样，才去单独学习。
          \item 1 部民法典 + 3 个司法解释（总则编、担保制度解释、民间借贷合同解释）。
        \end{enumerate}
  \item 民法典。7 编 1260 条。
        \begin{enumerate}
          \item 总则。10\%。8 分。
                \begin{enumerate}
                  \item 民事法律关系。民事主体（自然人、法人、非法人）。
                  \item 民事法律事实。民事法律行为、代理、诉讼时效。
                \end{enumerate}
          \item 财产权。
                \begin{enumerate}
                  \item 物权。所有权、用益物权、担保物权、准物权（占有）。10 分。
                  \item 合同。通则、19 类典型合同、准合同（无因管理 + 不当得利）。
                \end{enumerate}
          \item 人身权。5 分。
                \begin{enumerate}
                  \item 人格权。2 类一般人格权、9 类具体人格权、个人信息保护。
                  \item 婚姻家庭。结婚、家庭关系、离婚、收养。
                \end{enumerate}
          \item 人身 + 财产（综合性权利）。
                \begin{enumerate}
                  \item 继承。法定继承 vs 遗嘱继承、代位继承 vs 转继承、遗赠 vs 遗赠抚养协议。
                \end{enumerate}
          \item 侵权责任。共同侵权、精神损害赔偿、7 类特殊主体责任、7 类典型侵权责任。归责原则、免责事由。
        \end{enumerate}
\end{enumerate}
